% =====================================================================
\chapter{Información y entropía}
\label{tema:1}
% =====================================================================

Este tema define la medida de la información introducida por Shannon~\cite{shannon} y sus versiones para varias variables: entropía, entropía conjunta y entropía condicionada. Se sigue~\cite[§2.1, §2.2 y §2.5]{ct}; la interpretación de la entropía como número de preguntas y la discusión del contenido de información proceden de~\cite[cap.~2 y~4]{mackay} y de la primera clase de~\cite{camit}. En todo el tema, $X$ es una variable aleatoria discreta con valores en un alfabeto finito $\X$ y función de masa $p(x)=P(X=x)$.

\section{El contenido de información de un suceso}

Un suceso seguro no informa de nada; un suceso improbable, al ocurrir, informa mucho. Se busca una función $h$ que asigne a cada probabilidad $p\in(0,1]$ una cantidad de información y que cumpla tres requisitos naturales:
\begin{enumerate}
  \item $h(1)=0$: lo que es seguro no aporta información;
  \item $h$ es decreciente: cuanto menos probable es un suceso, más informa;
  \item $h(pq)=h(p)+h(q)$: la información de dos sucesos independientes es la suma de las informaciones.
\end{enumerate}

\begin{proposicion}[Caracterización del contenido de información]\label{prop:caracterizacion}
Sea $h\colon(0,1]\to\R$ decreciente (en sentido amplio) y tal que $h(pq)=h(p)+h(q)$ para todo $p,q\in(0,1]$. Entonces existe una constante $c\ge 0$ tal que $h(p)=-c\log_2 p$ para todo $p$.
\end{proposicion}

La constante $c$ solo fija la unidad de medida. Se adopta $c=1$.

\begin{definicion}[Contenido de información]\label{def:sorpresa}
El \emph{contenido de información} o \emph{sorpresa} del resultado $x$, con $p(x)>0$, es
\[
h(x)=-\log_2 p(x)\quad\text{bits}.
\]
Con el logaritmo natural la unidad es el \emph{nat}; como $\ln p=\ln 2\cdot\log_2 p$, se tiene $1\text{ nat}=1/\ln 2\approx 1{,}443$ bits.
\end{definicion}

\begin{center}
\begin{tabular}{@{}lcc@{}}
\toprule
Suceso & $p$ & $h$ (bits) \\
\midrule
cara en una moneda equilibrada & $1/2$ & $1$ \\
un~3 en un dado equilibrado & $1/6$ & $2{,}585$ \\
uno entre 8 resultados equiprobables & $1/8$ & $3$ \\
un token al que el modelo asigna $0{,}9$ & $0{,}9$ & $0{,}152$ \\
un token al que el modelo asigna $0{,}1$ & $0{,}1$ & $3{,}322$ \\
uno entre $1024$ resultados equiprobables & $2^{-10}$ & $10$ \\
\bottomrule
\end{tabular}
\end{center}

Si $p(x)=2^{-k}$, el resultado $x$ puede identificarse con $k$ preguntas de respuesta sí o no que dividan por la mitad, en cada paso, el conjunto de candidatos. El contenido de información generaliza esta cuenta a probabilidades arbitrarias~\cite[cap.~4]{mackay}.

\addtocounter{definicion}{1}%

\section{Entropía}

\begin{definicion}[Entropía]\label{def:entropia}
La \emph{entropía} de $X$ es la sorpresa esperada:
\[
H(X)=H(p)=\E\bigl[-\log p(X)\bigr]=-\sum_{x\in\X}p(x)\log p(x),
\]
con el convenio $0\log 0=0$, justificado por $\lim_{t\to 0^{+}}t\log t=0$.
\end{definicion}

La entropía depende solo de las probabilidades, no de los valores que toma $X$. Mide la incertidumbre sobre $X$ antes de observarla o, de forma equivalente, la información que se obtiene, en media, al observarla.

\addtocounter{definicion}{1}%

\addtocounter{definicion}{1}%

\addtocounter{figure}{1}%

\begin{proposicion}[Cotas elementales]\label{prop:cotas-entropia}
Para toda distribución $p$ sobre un alfabeto de $K$ símbolos,
\begin{enumerate}
  \item $H(p)\ge 0$, con igualdad si y solo si $p$ es degenerada, es decir, $p(x_0)=1$ para algún $x_0$;
  \item $H(p)\le\log K$, con igualdad si y solo si $p$ es la distribución uniforme.
\end{enumerate}
\end{proposicion}

\begin{center}
\begin{tabular}{@{}lc@{}}
\toprule
Distribución & Entropía (bits) \\
\midrule
uniforme sobre 27 símbolos (26 letras y el espacio) & $\log 27=4{,}755$ \\
$(1/2,1/4,1/4)$ & $1{,}5$ \\
moneda con $P(\text{cara})=0{,}1$ & $0{,}469$ \\
distribución degenerada & $0$ \\
\bottomrule
\end{tabular}
\end{center}

\section{Entropía conjunta y entropía condicionada}

Sean $X$ e $Y$ variables discretas con función de masa conjunta $p(x,y)$, marginales $p(x)$ y $p(y)$ y condicionada $p(y\mid x)=p(x,y)/p(x)$ cuando $p(x)>0$.

\begin{definicion}[Entropía conjunta y condicionada]\label{def:conjunta}
La \emph{entropía conjunta} de $(X,Y)$ y la \emph{entropía de $Y$ condicionada a $X$} son
\begin{gather*}
H(X,Y)=-\sum_{x,y}p(x,y)\log p(x,y),\\
H(Y\mid X)=\sum_{x}p(x)\,H(Y\mid X=x)=-\sum_{x,y}p(x,y)\log p(y\mid x),
\end{gather*}
donde $H(Y\mid X=x)=-\sum_y p(y\mid x)\log p(y\mid x)$ es la entropía de la distribución condicionada.
\end{definicion}

La entropía condicionada es la incertidumbre que queda sobre $Y$ cuando se conoce $X$, promediada sobre los valores de $X$.

\begin{teorema}[Regla de la cadena para la entropía]\label{teo:cadena-entropia}
$H(X,Y)=H(X)+H(Y\mid X)$.
\end{teorema}

\begin{corolario}[Regla de la cadena para secuencias]\label{cor:cadena-secuencias}
Para variables $X_1,\dots,X_n$,
\[
H(X_1,\dots,X_n)=\sum_{t=1}^{n}H(X_t\mid X_1,\dots,X_{t-1}).
\]
\end{corolario}

El \cref{cor:cadena-secuencias} es la versión informacional de la factorización autorregresiva del módulo~1, $P(x_{1:T})=\prod_t P(x_t\mid x_{1:t-1})$: tomando $-\log$ y esperanzas en esa factorización se obtiene exactamente esta suma. La incertidumbre de un texto completo es la suma de las incertidumbres de cada token dado su contexto.

\addtocounter{definicion}{1}%

\begin{observacion}
En general $H(Y\mid X)\ne H(X\mid Y)$, pero $H(X)-H(X\mid Y)=H(Y)-H(Y\mid X)$, porque ambas diferencias valen $H(X)+H(Y)-H(X,Y)$ por la regla de la cadena aplicada en los dos órdenes. Esta cantidad común es la información mutua del \cref{tema:2}.
\end{observacion}

\section{Tasa de entropía y redundancia del lenguaje}

\begin{definicion}[Tasa de entropía]\label{def:tasa}
La \emph{tasa de entropía} de un proceso $X_1,X_2,\dots$ es
\[
\mathcal{H}=\lim_{n\to\infty}\frac1n H(X_1,\dots,X_n),
\]
cuando el límite existe. Se mide en bits por símbolo.
\end{definicion}

Si las $X_t$ son independientes e idénticamente distribuidas, $H(X_1,\dots,X_n)=nH(X_1)$ y $\mathcal{H}=H(X_1)$. En un texto, en cambio, los símbolos dependen fuertemente del contexto y $\mathcal{H}$ es mucho menor que la entropía de un símbolo aislado~\cite[cap.~4]{ct}. Shannon~\cite{shannon51} estimó la tasa de entropía del inglés pidiendo a personas que adivinaran un texto letra a letra, y obtuvo entre $0{,}6$ y $1{,}3$ bits por letra, muy por debajo de los $\log 27=4{,}755$ bits de un alfabeto de 27 símbolos equiprobables. La diferencia mide la \emph{redundancia} del lenguaje: la parte del texto que es previsible a partir de lo anterior.

\addtocounter{definicion}{1}%

